On the same seven-input strong source, the literal native isospin raising
operator transports the complete neutron ground-ray certificate to the
proton ray. Their mass difference is $0.00129333251$ GeV. The intrinsic
neutron charge is positive at graph radius $1/3$ and negative at $2/3$.
A maximally mixed input on all $120$ physical colour singlets gives a
ground-ray filter probability $1/120$; all $2024$ occupation effects and
their failure records are retained through a compressed positive-operator
measurement. Its nonselective mean source work on that input is zero.

With the shared quadrupole/string length near $0.7559087287$ fm, the
orientation-averaged charge records give a proton root-mean-square radius
near $0.3345681563$ fm and neutron charge second moment near
$-0.03731193543$ fm$^2$. The common length cancels in
$-\langle r_n^2\rangle/\langle r_p^2\rangle$. Its certified interval is
disjoint from the comparison box formed by the CODATA2022 proton radius
and PDG2025 neutron mean-square radius. No single length therefore fits
both radii on this selected source. The graph orientation average is
specified; no relativistic Sachs-current or continuum nucleon-pole
identification follows from these records
(\texttt{model/openwave\_baryon\_form\_factors}).
